\chapter{Beschleuniger}

\section{Ziel dieses Kapitels}

Dieses Kapitel behandelt Teilchenbeschleuniger.
Beschleuniger sind zentrale Werkzeuge der Kern- und Teilchenphysik, weil sie
geladene Teilchen auf hohe Energien bringen und damit Streuung, Kernreaktionen
oder Teilchenkollisionen ermoeglichen.

\begin{notebox}
Dieses Kapitel behandelt die Stichwoerter:
\begin{itemize}
    \item Beschleuniger,
    \item elektrische Beschleunigung,
    \item Radiofrequenz-Kavitaeten,
    \item Linearbeschleuniger,
    \item Kreisbeschleuniger,
    \item Zyklotron,
    \item Synchrotron,
    \item Ablenkmagnete,
    \item Quadrupolfokussierung,
    \item Collider,
    \item Fixed-Target-Experiment,
    \item LHC,
    \item Strahlenergie,
    \item Luminositaet.
\end{itemize}
\end{notebox}

\section{Warum braucht man Beschleuniger?}

In der Kern- und Teilchenphysik will man sehr kleine Strukturen untersuchen.
Dafuer braucht man Teilchen mit hoher Energie und kleinem de-Broglie-Wellenlaengenmassstab.

\begin{conceptbox}
Je groesser der Impuls eines Teilchens ist, desto kleiner ist seine de-Broglie-Wellenlaenge:
\[
    \lambda = \frac{h}{p}.
\]

Kleine Wellenlaengen erlauben die Untersuchung kleiner Strukturen.
Beschleuniger liefern daher die notwendige Energie und den notwendigen Impuls.
\end{conceptbox}

\begin{conceptbox}
Beschleuniger werden verwendet fuer:
\begin{itemize}
    \item Streuexperimente,
    \item Kernreaktionen,
    \item Erzeugung neuer Teilchen,
    \item Untersuchung von Formfaktoren,
    \item Teilchenphysik bei hohen Schwerpunktsenergien,
    \item medizinische und technische Anwendungen.
\end{itemize}
\end{conceptbox}

\begin{examtipbox}
Merksatz:
Hohe Energie bedeutet kurze Wellenlaenge.
Kurze Wellenlaenge bedeutet hohe Strukturaufloesung.
\end{examtipbox}

\section{Welche Teilchen kann man beschleunigen?}

Teilchenbeschleuniger beschleunigen geladene Teilchen.
Neutrale Teilchen koennen nicht direkt durch elektrische Felder beschleunigt werden.

\begin{definitionbox}
Ein \emph{Teilchenbeschleuniger} ist eine Anlage, die geladene Teilchen durch
elektrische Felder auf hohe Energien bringt.
\end{definitionbox}

\begin{conceptbox}
Beschleunigt werden koennen zum Beispiel:
\begin{itemize}
    \item Elektronen,
    \item Positronen,
    \item Protonen,
    \item Ionen,
    \item schwere Kerne.
\end{itemize}

Neutronen sind elektrisch neutral und koennen daher nicht direkt durch elektrische
Felder beschleunigt werden.
\end{conceptbox}

\begin{examtipbox}
Elektrische Felder beschleunigen geladene Teilchen.

Magnetische Felder lenken geladene Teilchen ab, verrichten aber keine Arbeit am
Teilchen.
\end{examtipbox}

\section{Lorentzkraft}

Die Bewegung geladener Teilchen in elektrischen und magnetischen Feldern wird durch
die Lorentzkraft beschrieben.

\begin{formulabox}
Die Lorentzkraft ist
\[
    \vec F
    =
    q\left(
        \vec E
        +
        \vec v \times \vec B
    \right).
\]
Dabei ist:
\begin{itemize}
    \item $q$ die Teilchenladung,
    \item $\vec E$ das elektrische Feld,
    \item $\vec B$ das magnetische Feld,
    \item $\vec v$ die Teilchengeschwindigkeit.
\end{itemize}
\end{formulabox}

\begin{conceptbox}
Das elektrische Feld kann Arbeit verrichten:
\[
    W = qU.
\]
Dadurch steigt die kinetische Energie des Teilchens.

Das magnetische Feld wirkt senkrecht zur Geschwindigkeit.
Es aendert daher die Bewegungsrichtung, aber nicht direkt die Energie.
\end{conceptbox}

\begin{examtipbox}
Beschleunigung:
\[
    \vec E \text{-Feld}.
\]

Ablenkung und Kreisbahn:
\[
    \vec B \text{-Feld}.
\]
\end{examtipbox}

\section{Energiegewinn in einem elektrischen Potential}

Ein geladenes Teilchen gewinnt Energie, wenn es eine elektrische Spannung durchlaeuft.

\begin{formulabox}
Der Energiegewinn ist
\[
    \Delta E = qU.
\]
Fuer ein Teilchen mit Ladung $q=ze$ gilt:
\[
    \Delta E = zeU.
\]
\end{formulabox}

\begin{definitionbox}
Ein Elektronenvolt ist die Energie, die ein Teilchen mit Elementarladung $e$ beim
Durchlaufen einer Spannung von einem Volt gewinnt:
\[
    1\,\mathrm{eV}=e\cdot 1\,\mathrm{V}.
\]
\end{definitionbox}

\begin{examplebox}
Ein Proton durchlaeuft eine Spannung von
\[
    U=\SI{1}{MV}.
\]
Da das Proton die Ladung
\[
    q=e
\]
hat, gewinnt es die Energie
\[
    \Delta E=eU=\SI{1}{MeV}.
\]
\end{examplebox}

\begin{conceptbox}
Um ein Proton auf
\[
    \SI{7}{TeV}
\]
zu bringen, braucht man insgesamt einen Energiegewinn von
\[
    \SI{7}{TeV}.
\]
Das entspricht fuer ein einfach geladenes Teilchen einer Gesamtspannung von
\[
    \SI{7}{TV}.
\]
In einem Kreisbeschleuniger wird diese Spannung nicht auf einmal durchlaufen,
sondern schrittweise bei vielen Umlaeufen aufgebaut.
\end{conceptbox}

\section{Statische Beschleuniger}

Die einfachste Beschleunigung erfolgt durch eine statische Hochspannung.

\begin{definitionbox}
Ein \emph{statischer Beschleuniger} beschleunigt geladene Teilchen durch eine feste
Potentialdifferenz.
\end{definitionbox}

\begin{conceptbox}
Beispiele fuer statische Beschleuniger sind:
\begin{itemize}
    \item Cockcroft-Walton-Beschleuniger,
    \item Van-de-Graaff-Beschleuniger.
\end{itemize}

Die erreichbare Energie ist durch die maximal technisch beherrschbare Spannung
begrenzt.
\end{conceptbox}

\begin{definitionbox}
Ein \emph{Van-de-Graaff-Beschleuniger} erzeugt eine hohe elektrische Spannung durch
Transport von Ladung auf eine isolierte Elektrode.
\end{definitionbox}

\begin{examtipbox}
Statische Beschleuniger sind konzeptionell einfach:
\[
    \Delta E=qU.
\]
Ihr Nachteil ist, dass sehr hohe Spannungen technisch schwer zu realisieren sind.
\end{examtipbox}

\section{Radiofrequenz-Kavitaeten}

Moderne Hochenergie-Beschleuniger verwenden zeitabhaengige elektrische Felder in
Radiofrequenz-Kavitaeten.

\begin{definitionbox}
Eine \emph{Radiofrequenz-Kavitaet} beziehungsweise \emph{RF-Kavitaet} ist ein
Hohlraumresonator, in dem ein oszillierendes elektrisches Feld aufgebaut wird.
Dieses Feld beschleunigt geladene Teilchen.
\end{definitionbox}

\begin{conceptbox}
Die Beschleunigung erfolgt nur dann effizient, wenn die Teilchen zum richtigen
Zeitpunkt in der Kavitaet ankommen.

Das elektrische Feld muss so phasenrichtig sein, dass es das Teilchen in
Bewegungsrichtung beschleunigt.
\end{conceptbox}

\begin{notebox}
Die LHC-Folien beschreiben die Beschleunigung mit Hochfrequenzresonatoren.
In einem Kreisbeschleuniger erfolgt die Beschleunigung schrittweise:
Runde fuer Runde erhalten die Teilchen einen kleinen Energiezuwachs.
\end{notebox}

\begin{examtipbox}
RF-Kavitaeten liefern die Energieerhoehung.

Magnete halten und fokussieren den Strahl.
\end{examtipbox}

\section{Linearbeschleuniger}

Bei einem Linearbeschleuniger bewegen sich die Teilchen entlang einer geraden Linie.

\begin{definitionbox}
Ein \emph{Linearbeschleuniger} beschleunigt geladene Teilchen entlang einer geraden
Bahn durch eine Folge elektrischer Beschleunigungsstrecken.
\end{definitionbox}

\begin{conceptbox}
Ein Linearbeschleuniger hat den Vorteil, dass keine Synchrotronstrahlung durch
Kreisbewegung entsteht.

Der Nachteil ist:
Jede Beschleunigungsstruktur wird im Wesentlichen nur einmal durchlaufen.
Fuer sehr hohe Energien muss der Beschleuniger daher sehr lang sein.
\end{conceptbox}

\begin{notebox}
Linearbeschleuniger werden oft auch als Vorbeschleuniger verwendet.
Teilchen werden zuerst in einem Linearbeschleuniger auf eine Anfangsenergie gebracht
und danach in einen Kreisbeschleuniger eingespeist.
\end{notebox}

\begin{examtipbox}
Linearbeschleuniger:
\begin{itemize}
    \item gerade Bahn,
    \item keine Umlaeufe,
    \item kein Energieverlust durch Kreisbewegung,
    \item sehr lang bei hohen Energien.
\end{itemize}
\end{examtipbox}

\section{Kreisbeschleuniger}

Bei einem Kreisbeschleuniger laufen Teilchen immer wieder durch dieselben
Beschleunigungsstrecken.

\begin{definitionbox}
Ein \emph{Kreisbeschleuniger} fuehrt geladene Teilchen auf einer geschlossenen
Bahn und beschleunigt sie bei vielen Umlaeufen schrittweise.
\end{definitionbox}

\begin{conceptbox}
Der Vorteil eines Kreisbeschleunigers ist:
Dieselben RF-Kavitaeten koennen viele Male verwendet werden.

Dadurch kann man hohe Endenergien erreichen, ohne eine entsprechend lange gerade
Beschleunigungsstrecke bauen zu muessen.
\end{conceptbox}

\begin{conceptbox}
Der Nachteil ist:
Geladene Teilchen auf Kreisbahnen werden beschleunigt, weil sich ihre Richtung
aendert.
Dadurch koennen sie Synchrotronstrahlung aussenden.

Dieser Effekt ist besonders stark fuer leichte Teilchen wie Elektronen.
\end{conceptbox}

\section{Magnetische Ablenkung}

Damit ein Teilchen eine Kreisbahn beschreibt, braucht man ein Magnetfeld.

\begin{formulabox}
Fuer ein Teilchen mit Impuls $p$ und Ladung $q$ in einem Magnetfeld $B$ gilt
naeherungsweise:
\[
    p = qBr.
\]
Fuer hochrelativistische Teilchen schreibt man oft:
\[
    p[\mathrm{GeV}/c]
    \approx
    0.3\,B[\mathrm{T}]\,r[\mathrm{m}].
\]
\end{formulabox}

\begin{conceptbox}
Aus
\[
    p=qBr
\]
folgt:
Je groesser der Impuls $p$ ist, desto staerker muss das Magnetfeld $B$ sein
oder desto groesser muss der Radius $r$ des Beschleunigers sein.

Das erklaert, warum Hochenergie-Kreisbeschleuniger sehr grosse Ringe und starke
Magnete benoetigen.
\end{conceptbox}

\begin{examtipbox}
Die maximale Energie in einem Kreisbeschleuniger wird wesentlich begrenzt durch:
\begin{itemize}
    \item Magnetfeldstaerke,
    \item Ringradius,
    \item Synchrotronstrahlung,
    \item Strahlstabilitaet.
\end{itemize}
\end{examtipbox}

\section{Zyklotron}

Das Zyklotron ist ein klassischer Kreisbeschleuniger fuer geladene Teilchen.

\begin{definitionbox}
Ein \emph{Zyklotron} beschleunigt geladene Teilchen in einem konstanten Magnetfeld
und einem hochfrequenten elektrischen Wechselfeld.
\end{definitionbox}

\begin{conceptbox}
Im Zyklotron bewegen sich die Teilchen auf Spiralbahnen nach aussen.
Bei jedem Durchgang durch den Spalt zwischen den Elektroden erhalten sie Energie.

Fuer nichtrelativistische Teilchen ist die Umlauffrequenz unabhaengig vom Radius:
\[
    \omega = \frac{qB}{m}.
\]
\end{conceptbox}

\begin{formulabox}
Die Zyklotronfrequenz ist:
\[
    \omega_c = \frac{qB}{m}.
\]
Die Frequenz ist:
\[
    f_c = \frac{qB}{2\pi m}.
\]
\end{formulabox}

\begin{conceptbox}
Bei hohen Energien wird die relativistische Massenzunahme beziehungsweise der
Lorentzfaktor wichtig.
Dann aendert sich die Umlauffrequenz, und das einfache Zyklotron geraet ausser Takt.
\end{conceptbox}

\begin{examtipbox}
Zyklotron:
\begin{itemize}
    \item konstantes Magnetfeld,
    \item Hochfrequenzfeld zur Beschleunigung,
    \item Spiralbahn,
    \item gut fuer nichtrelativistische Teilchen,
    \item relativistische Effekte begrenzen die einfache Bauform.
\end{itemize}
\end{examtipbox}

\section{Synchrotron}

Das Synchrotron ist ein moderner Kreisbeschleuniger, bei dem Magnetfelder und
RF-Frequenzen an die steigende Teilchenenergie angepasst werden.

\begin{definitionbox}
Ein \emph{Synchrotron} ist ein Kreisbeschleuniger, bei dem die Teilchen auf einer
nahezu festen Bahn umlaufen und Magnetfeld sowie Beschleunigungsfrequenz synchron
zur Teilchenenergie angepasst werden.
\end{definitionbox}

\begin{conceptbox}
Im Synchrotron bleibt der Bahnradius im Wesentlichen konstant.

Wenn die Teilchenenergie und damit der Impuls steigen, muss das Magnetfeld entsprechend
erhoeht werden:
\[
    p=qBr.
\]
\end{conceptbox}

\begin{notebox}
Die Vorlesungsabbildungen zu Collider und Detektoren enthalten ein Elektronen-Synchrotron
und einen Ringaufbau mit:
\begin{itemize}
    \item Ablenkmagneten,
    \item Quadrupolfokussiereinheiten,
    \item Beschleunigungseinheiten.
\end{itemize}
\end{notebox}

\begin{examtipbox}
Synchrotron:
\begin{itemize}
    \item feste Kreisbahn,
    \item Magnetfeld wird mit der Energie hochgefahren,
    \item RF-System bleibt phasensynchron,
    \item geeignet fuer sehr hohe Energien.
\end{itemize}
\end{examtipbox}

\section{Synchrotronstrahlung}

Geladene Teilchen, die auf Kreisbahnen gezwungen werden, strahlen elektromagnetische
Strahlung ab.

\begin{definitionbox}
\emph{Synchrotronstrahlung} ist elektromagnetische Strahlung, die von geladenen
Teilchen bei beschleunigter Bewegung auf gekruemmten Bahnen ausgesendet wird.
\end{definitionbox}

\begin{conceptbox}
Synchrotronstrahlung ist besonders stark fuer leichte Teilchen.
Elektronen verlieren in Kreisbeschleunigern viel Energie durch Synchrotronstrahlung.

Protonen sind viel schwerer.
Deshalb ist Synchrotronstrahlung bei Protonen bei gleicher Energie deutlich weniger
problematisch.
\end{conceptbox}

\begin{examtipbox}
Warum sind sehr hohe Energien fuer Elektronen in Kreisbeschleunigern schwierig?

Weil Elektronen wegen ihrer kleinen Masse starke Synchrotronstrahlung aussenden.
Daher sind fuer sehr hohe Elektronenergien Linearbeschleuniger attraktiv.
\end{examtipbox}

\section{Fokussierung des Strahls}

Ein Teilchenstrahl muss nicht nur beschleunigt, sondern auch gebuendelt werden.
Sonst wuerde er auseinanderlaufen.

\begin{definitionbox}
\emph{Strahlfokussierung} bedeutet, den Teilchenstrahl raeumlich zusammenzuhalten
und auf kleine Querschnitte zu buendeln.
\end{definitionbox}

\begin{conceptbox}
In Beschleunigern verwendet man Magnetoptik.
Dipolmagnete lenken den Strahl.
Quadrupolmagnete fokussieren den Strahl.

Ein Quadrupol fokussiert in einer Ebene und defokussiert in der anderen.
Durch geeignete Kombination mehrerer Quadrupole kann der Strahl insgesamt stabil
gefuehrt werden.
\end{conceptbox}

\begin{definitionbox}
Ein \emph{Ablenkmagnet} beziehungsweise Dipolmagnet kraemmt die Bahn des Teilchenstrahls.

Ein \emph{Quadrupolmagnet} fokussiert den Teilchenstrahl.
\end{definitionbox}

\begin{examtipbox}
Dipolmagnete:
Bahn kruemmen.

Quadrupolmagnete:
Strahl fokussieren.
\end{examtipbox}

\section{Collider und Fixed-Target-Experimente}

Es gibt zwei wichtige experimentelle Geometrien:
Fixed-Target-Experimente und Collider-Experimente.

\begin{definitionbox}
Bei einem \emph{Fixed-Target-Experiment} trifft ein beschleunigter Strahl auf ein
ruhendes Target.

Bei einem \emph{Collider} kollidieren zwei gegenlaeufige Teilchenstrahlen.
\end{definitionbox}

\begin{conceptbox}
Der Vorteil eines Colliders ist:
Ein viel groesserer Anteil der Strahlenergie steht im Schwerpunktsystem fuer die
Erzeugung neuer Teilchen zur Verfuegung.

Bei einem Fixed-Target-Experiment steckt ein grosser Teil der Energie nach der
Kollision in der Bewegung des Schwerpunktsystems.
\end{conceptbox}

\begin{formulabox}
Fuer zwei gleich energiereiche, frontal kollidierende Teilchenstrahlen gilt:
\[
    E_{\mathrm{CM}}
    \approx
    E_1+E_2.
\]
Fuer einen Collider mit zwei Protonstrahlen von je
\[
    \SI{7}{TeV}
\]
gilt daher:
\[
    E_{\mathrm{CM}}=\SI{14}{TeV}.
\]
\end{formulabox}

\begin{examtipbox}
Collider sind fuer Teilchenphysik sehr effizient, weil sie hohe
Schwerpunktsenergien erreichen.

Fixed-Target-Experimente koennen trotzdem sinnvoll sein, zum Beispiel fuer hohe
Intensitaeten, spezielle Targets oder Praezisionsmessungen.
\end{examtipbox}

\section{Der LHC}

Der Large Hadron Collider, kurz LHC, ist ein Proton-Proton-Collider am CERN.

\begin{definitionbox}
Der \emph{Large Hadron Collider} ist ein Kreisbeschleuniger und Collider, in dem
hauptsaechlich Protonenstrahlen mit sehr hoher Energie zur Kollision gebracht werden.
\end{definitionbox}

\begin{notebox}
Die LHC-Folien betonen:
\begin{itemize}
    \item LHC steht fuer Large Hadron Collider.
    \item Das beschleunigte Hadron ist hauptsaechlich das Proton.
    \item Es koennen auch Kerne wie Blei oder Xenon beschleunigt werden.
    \item Der LHC ist ein Collider, nicht nur ein Fixed-Target-Beschleuniger.
\end{itemize}
\end{notebox}

\begin{conceptbox}
Im LHC laufen zwei Protonenstrahlen in entgegengesetzten Richtungen.
An bestimmten Kollisionspunkten werden sie zur Ueberlappung gebracht.

Dort befinden sich grosse Detektoren wie ATLAS, CMS, ALICE und LHCb.
\end{conceptbox}

\begin{notebox}
Die LHC-Detektorfolien nennen:
\begin{itemize}
    \item ATLAS = A Toroidal LHC ApparatuS,
    \item CMS = Compact Muon Solenoid,
    \item ALICE = A Large Ion Collider Experiment,
    \item LHCb = LHC beauty.
\end{itemize}
\end{notebox}

\section{Beschleunigung im LHC}

\begin{conceptbox}
Im LHC erfolgt die Beschleunigung mit Radiofrequenz-Kavitaeten.
Die Teilchen erhalten pro Umlauf einen kleinen Energiezuwachs.

Da die Protonen sehr viele Umlaeufe machen, summieren sich diese kleinen
Energiegewinne zu einer sehr grossen Endenergie.
\end{conceptbox}

\begin{notebox}
Die LHC-Folien geben als Orientierung:
\begin{itemize}
    \item Endenergie pro Proton etwa $\SI{7}{TeV}$,
    \item ein Beschleunigungsschritt pro Umlauf bringt grob $\SI{0.7}{MeV}$,
    \item der Umfang des Rings liegt bei etwa $\SI{30}{km}$,
    \item die Protonen bewegen sich praktisch mit Lichtgeschwindigkeit.
\end{itemize}
\end{notebox}

\begin{calculationbox}
Grobe Anzahl der Umlaeufe fuer die Beschleunigung von vernachlaessigbarer
Einschussenergie auf
\[
    \SI{7}{TeV}
\]
bei einem Energiegewinn von
\[
    \SI{0.7}{MeV}
\]
pro Umlauf:
\[
    n
    \approx
    \frac{\SI{7}{TeV}}{\SI{0.7}{MeV}}.
\]
Mit
\[
    \SI{7}{TeV}=\SI{7e6}{MeV}
\]
folgt:
\[
    n
    \approx
    \frac{7\cdot 10^6}{0.7}
    =
    10^7.
\]
Also sind grob zehn Millionen Umlaeufe noetig.
\end{calculationbox}

\section{Relativistische Geschwindigkeit im LHC}

Bei TeV-Energien sind Protonen hochrelativistisch.
Ihre Geschwindigkeit liegt extrem nahe bei der Lichtgeschwindigkeit.

\begin{formulabox}
Die Gesamtenergie eines Teilchens ist:
\[
    E = \gamma mc^2.
\]
Daraus:
\[
    \gamma = \frac{E}{mc^2}.
\]
Die Geschwindigkeit ist gegeben durch:
\[
    \beta=\frac{v}{c}
    =
    \sqrt{1-\frac{1}{\gamma^2}}.
\]
\end{formulabox}

\begin{conceptbox}
Wenn die Energie bereits sehr gross gegenueber der Ruheenergie ist, fuehrt eine
weitere Energieerhoehung kaum noch zu einer merklichen Erhoehung von $v$.

Stattdessen steigt vor allem der Impuls und die Gesamtenergie.
\end{conceptbox}

\begin{examtipbox}
Bei hochrelativistischen Teilchen gilt:
Die Geschwindigkeit ist fast $c$.
Mehr Energie bedeutet nicht viel mehr Geschwindigkeit, sondern vor allem mehr Impuls
und mehr Schwerpunktsenergie.
\end{examtipbox}

\section{Protonenpakete und Strahlstruktur}

Teilchenstrahlen in Beschleunigern sind nicht kontinuierlich homogen verteilt.
Sie bestehen aus Paketen.

\begin{definitionbox}
Ein \emph{Bunch} beziehungsweise Teilchenpaket ist eine Gruppe vieler Teilchen, die
gemeinsam im Beschleuniger umlaeuft.
\end{definitionbox}

\begin{conceptbox}
Im LHC sind die Protonen in Paketen organisiert.
Diese Pakete kollidieren an den Wechselwirkungspunkten.

Die Zahl der Kollisionen haengt davon ab:
\begin{itemize}
    \item wie viele Pakete im Ring sind,
    \item wie viele Protonen in jedem Paket sind,
    \item wie klein die Strahlen fokussiert werden,
    \item wie oft sich die Pakete treffen.
\end{itemize}
\end{conceptbox}

\begin{notebox}
Die LHC-Folien nennen als typische Groessen:
\begin{itemize}
    \item etwa $10^{11}$ Protonen pro Paket,
    \item bis zu $2808$ Protonenpakete pro Strahl,
    \item Protonenergie von etwa $\SI{7}{TeV}$.
\end{itemize}
\end{notebox}

\section{Gesamtenergie des Protonenstrahls}

Auch wenn ein einzelnes Proton mikroskopisch ist, enthaelt ein ganzer LHC-Strahl
makroskopisch grosse Energie.

\begin{calculationbox}
Verwende die Groessen:
\[
    E_p=\SI{7}{TeV},
    \qquad
    N=1.15\cdot 10^{11}
\]
Protonen pro Paket und
\[
    k=2808
\]
Pakete pro Strahl.

Die Gesamtenergie ist:
\[
    E_{\mathrm{Strahl}}
    =
    E_p\,N\,k.
\]

Mit
\[
    \SI{1}{eV}=1.602\cdot 10^{-19}\,\si{J}
\]
ist:
\[
    \SI{7}{TeV}
    =
    7\cdot 10^{12}\,\mathrm{eV}
    \cdot
    1.602\cdot 10^{-19}\,\si{J/eV}.
\]
Also:
\[
    \SI{7}{TeV}
    \approx
    1.12\cdot 10^{-6}\,\si{J}.
\]

Damit:
\[
    E_{\mathrm{Strahl}}
    \approx
    1.12\cdot 10^{-6}\,\si{J}
    \cdot
    1.15\cdot 10^{11}
    \cdot
    2808.
\]
\[
    E_{\mathrm{Strahl}}
    \approx
    3.6\cdot 10^{8}\,\si{J}
    =
    \SI{360}{MJ}.
\]
\end{calculationbox}

\begin{conceptbox}
Die LHC-Folien vergleichen diese Energie mit der kinetischen Energie eines schweren
Zuges.

Das zeigt:
Ein Teilchenstrahl ist nicht nur mikroskopisch interessant, sondern auch technisch
gefaehrlich.
Die Strahlkontrolle ist daher extrem wichtig.
\end{conceptbox}

\section{Luminositaet}

Bei Collidern ist nicht nur die Energie wichtig.
Man braucht auch viele Kollisionen.

\begin{definitionbox}
Die \emph{Luminositaet} $L$ beschreibt, wie viele Teilchenkollisionen pro Zeit und
Wirkungsquerschnitt moeglich sind.
\end{definitionbox}

\begin{formulabox}
Die Ereignisrate $R$ eines Prozesses ist:
\[
    R = L\sigma.
\]
Dabei ist:
\begin{itemize}
    \item $R$ die Ereignisrate,
    \item $L$ die Luminositaet,
    \item $\sigma$ der Wirkungsquerschnitt des Prozesses.
\end{itemize}
\end{formulabox}

\begin{conceptbox}
Ein seltener Prozess hat einen kleinen Wirkungsquerschnitt.
Um ihn trotzdem oft genug zu beobachten, braucht man eine hohe Luminositaet.

Deshalb muessen Collider nicht nur hohe Energie liefern, sondern auch sehr intensive
und gut fokussierte Strahlen.
\end{conceptbox}

\begin{examtipbox}
Energie entscheidet, welche Prozesse kinematisch moeglich sind.

Luminositaet entscheidet, wie oft diese Prozesse beobachtet werden.
\end{examtipbox}

\section{Was begrenzt die Energie eines Beschleunigers?}

\begin{conceptbox}
Die erreichbare Energie eines Beschleunigers wird durch mehrere Faktoren begrenzt:
\begin{itemize}
    \item maximale elektrische Felder in den RF-Kavitaeten,
    \item maximale Magnetfelder zur Strahlablenkung,
    \item Groesse beziehungsweise Radius des Rings,
    \item Synchrotronstrahlung,
    \item Strahlstabilitaet,
    \item Energie im Gesamtstrahl,
    \item technische und finanzielle Grenzen.
\end{itemize}
\end{conceptbox}

\begin{examtipbox}
Bei Kreisbeschleunigern ist besonders wichtig:
\[
    p=qBr.
\]
Hohe Teilchenimpulse brauchen entweder starke Magnete oder grosse Radien.
\end{examtipbox}

\section{Mini-Zusammenfassung}

\begin{notebox}
Die wichtigsten Punkte dieses Kapitels sind:
\begin{itemize}
    \item Beschleuniger bringen geladene Teilchen auf hohe Energien.
    \item Elektrische Felder beschleunigen.
    \item Magnetische Felder lenken und fokussieren.
    \item Der Energiegewinn in einer Spannung ist
    \[
        \Delta E=qU.
    \]
    \item RF-Kavitaeten beschleunigen Teilchen mit zeitabhaengigen elektrischen Feldern.
    \item Linearbeschleuniger beschleunigen entlang einer geraden Strecke.
    \item Kreisbeschleuniger nutzen dieselben Beschleunigungsstrecken viele Male.
    \item Im Zyklotron laufen Teilchen spiralfoermig nach aussen.
    \item Im Synchrotron bleibt die Bahn naeherungsweise konstant.
    \item Dipolmagnete lenken den Strahl.
    \item Quadrupolmagnete fokussieren den Strahl.
    \item Collider erreichen hohe Schwerpunktsenergien.
    \item Der LHC ist ein Proton-Proton-Collider.
    \item Luminositaet bestimmt zusammen mit dem Wirkungsquerschnitt die Ereignisrate.
\end{itemize}
\end{notebox}

\section{Pruefungscheck}

\begin{examtipbox}
Nach diesem Kapitel solltest du folgende Fragen beantworten koennen:
\begin{enumerate}
    \item Warum braucht man Beschleuniger in der Kern- und Teilchenphysik?
    \item Welche Teilchen kann man direkt beschleunigen?
    \item Welche Rolle spielt das elektrische Feld?
    \item Welche Rolle spielt das magnetische Feld?
    \item Wie lautet die Lorentzkraft?
    \item Was bedeutet $\Delta E=qU$?
    \item Was ist ein Van-de-Graaff-Beschleuniger?
    \item Was ist eine RF-Kavitaet?
    \item Wie funktioniert ein Linearbeschleuniger?
    \item Wie funktioniert ein Kreisbeschleuniger?
    \item Was ist ein Zyklotron?
    \item Warum funktioniert das einfache Zyklotron bei hohen Energien nicht mehr ideal?
    \item Was ist ein Synchrotron?
    \item Was ist Synchrotronstrahlung?
    \item Warum ist Synchrotronstrahlung besonders fuer Elektronen wichtig?
    \item Was machen Dipolmagnete?
    \item Was machen Quadrupolmagnete?
    \item Was ist der Unterschied zwischen Collider und Fixed-Target-Experiment?
    \item Warum liefert ein Collider hoehere Schwerpunktsenergie?
    \item Was ist der LHC?
    \item Warum sind LHC-Protonen trotz riesiger Energie kaum schneller als Licht?
    \item Was ist Luminositaet?
    \item Wie haengt die Ereignisrate mit Luminositaet und Wirkungsquerschnitt zusammen?
    \item Was begrenzt die maximale Energie eines Kreisbeschleunigers?
\end{enumerate}
\end{examtipbox}